\section{The arithmetic complex and its orbit bound}
\label{ac-section}

This section proves the geometric input in
\eqref{ht-geometric-input}.  In particular, the homology estimate in
Proposition~\ref{ht-arithmetic-numerical} no longer requires that input
as an additional hypothesis.  The construction is the ordered nerve
of the normal cones of the Hermitian reduction polyhedron used in
the submitted source\cite{alpo-uniform-source}.  We give the exact
arithmetic definitions and prove its exactness, the stabilizer bound,
and its orbit count.  This does not prove the separate global Selmer
class count discussed at the end of Section~\ref{ht-section}.

\subsection{The lattice and the Hermitian reduction polyhedron}

Let $F$ be a totally imaginary number field, let
$n=[F:\mathbb Q]=2s$, and let
$D=|\operatorname{disc}F|$.  Fix one embedding
$\sigma_v:F\longrightarrow\mathbb C$ from each conjugate pair.
Let $N\ge1$, and put
\[
 \Lambda=\mathcal O_F^N,\qquad
 d=nN,\qquad
 \Gamma=\operatorname{GL}_N(\mathcal O_F),\qquad
 \mathcal H=\prod_{v=1}^s\operatorname{Herm}_N(\mathbb C).
\]
Here $\operatorname{Herm}_N(\mathbb C)$ is the real vector space of
Hermitian $N$ by $N$ matrices.  Its dimension is $N^2$, and
$\dim_{\mathbb R}\mathcal H=sN^2\le dN$.  For $H=(H_v)\in\mathcal H$
and $y\in\Lambda$, define
\[
 q_H(y)=\sum_{v=1}^s \sigma_v(y)^*H_v\sigma_v(y).
\]
Every summand is real.  The star denotes conjugate transpose;
no involution of the number field is assumed.
Write
\[
 \begin{aligned}
 \mathcal P&=\{H:H_v\text{ is positive semidefinite for every }v\},\\
 \mathcal R&=\{H:q_H(y)\ge1\text{ for every }0\ne y\in\Lambda\}.
 \end{aligned}
 \tag{AC1}
\]
The group acts by
\[
 (gH)_v=\sigma_v(g)^{-*}H_v\sigma_v(g)^{-1}.
 \tag{AC2}
\]
Consequently $q_{gH}(y)=q_H(g^{-1}y)$, and the action preserves
$\mathcal R$ and $\mathcal P$.

We use the elementary facts that $\mathcal O_F$ is a free abelian
group of rank $n$, that the field norm of a nonzero algebraic integer
is a nonzero integer, and that the discriminant of an integral basis
is nonzero.  The volume normalization needed below follows directly
from these facts.  If $\omega_1,\ldots,\omega_n$ is an integral basis,
the determinant of the matrix of all $n$ embeddings has absolute
value $\sqrt D$.  Replacing each conjugate pair of rows by its real
and imaginary parts divides its absolute determinant by $2$.
Thus the image of $\mathcal O_F$ in $\mathbb C^s=\mathbb R^n$
is a full lattice of covolume $2^{-s}\sqrt D$.

\begin{lemma}\label{ac-positive-gram}
Let $G$ be a real symmetric $r$ by $r$ matrix, with $r\ge1$.
If
\[
                 z^{\mathsf T}Gz\ge1
       \quad(0\ne z\in\mathbb Z^r),
\]
then $G$ is positive definite and $\det G\ge r^{-r}$.
\end{lemma}
\begin{proof}
First $G$ is positive semidefinite.  Otherwise continuity supplies
a rational vector on which its quadratic form is negative, and
clearing denominators contradicts the hypothesis.

We recall and prove the elementary lattice estimate used here.
If a measurable set $A\subset\mathbb R^r$ has volume greater than
the covolume of a lattice $L$, then two distinct points of $A$
differ by a lattice vector.  Indeed, let $P$ be a half-open
fundamental parallelepiped.  Integration over $P$ of
$\sum_{\ell\in L}\mathbf1_A(x+\ell)$ gives $\operatorname{vol}(A)$.
If no two such points existed, this integer-valued sum would be at
most one everywhere, contradicting the volume inequality.
Applying this assertion to one half of an open centrally symmetric
convex set proves that such a set of volume greater than
$2^r\operatorname{covol}(L)$ contains a nonzero lattice point:
the difference of the two points belongs to the original set.

Suppose first that $G$ is positive definite, and choose a real
invertible matrix $T$ with $G=T^{\mathsf T}T$.  Apply the preceding
argument to $T\mathbb Z^r$ and the cube $(-R,R)^r$, where
$R>|\det T|^{1/r}$.  For some $0\ne z\in\mathbb Z^r$,
\[
                  1\le\|Tz\|^2\le rR^2.
\]
Letting $R$ decrease to $|\det T|^{1/r}$ gives
$1\le r(\det G)^{1/r}$.  For a merely positive semidefinite $G$,
apply this bound to $G+\varepsilon I$ and let
$\varepsilon\downarrow0$.  Its determinant is at least $r^{-r}$
in the limit.  Hence $\det G>0$, and positive semidefiniteness
then implies positive definiteness.
\end{proof}

\begin{proposition}\label{ac-polyhedron}
The set $\mathcal R$ is nonempty, closed and convex, and every
$H\in\mathcal R$ is positive definite at every complex place.
The defining inequalities are locally finite on $\mathcal R$:
for every $H\in\mathcal R$, there are a neighborhood $U$ of $H$
and a finite set $T\subset\Lambda\setminus\{0\}$ such that
\[
          q_K(y)>1\qquad(K\in U,\ y\notin T).
          \tag{AC3}
\]
\end{proposition}
\begin{proof}
Closedness and convexity follow because each $q_H(y)$ is real
linear in $H$.  The tuple of identity matrices belongs to
$\mathcal R$.  Indeed, a nonzero coordinate $a\in\mathcal O_F$
of $y$ has $|\sigma_v(a)|\ge1$ for at least one $v$;
otherwise $|\operatorname{Norm}_{F/\mathbb Q}(a)|<1$.

Choose the integral basis above in every coordinate of $\Lambda$.
Its embedding identifies $\mathbb Z^d$ with a full lattice in
$\prod_v\mathbb C^N$, and identifies the associated real vector
spaces.  The real quadratic form $q_H$ therefore has a symmetric
Gram matrix $G_H$ in these coordinates.  For $H\in\mathcal R$,
Lemma~\ref{ac-positive-gram} makes $G_H$ positive definite.
Since the embedding is a real vector-space isomorphism, a vector
supported at any one complex place shows that $H_v$ is positive
definite at that place.

Finally choose $c>0$ and a sufficiently small neighborhood $U$
such that
\[
 q_K(z)\ge c\sum_v\|z_v\|^2
       \quad(K\in U,\ z\in\prod_v\mathbb C^N).
\]
The lattice has only finitely many points in the ball
$\sum_v\|\sigma_v(y)\|^2\le c^{-1}$.  Taking these points for $T$
proves (AC3).
\end{proof}

For $H\in\mathcal R$ let
\[
                  M(H)=\{y\in\Lambda:q_H(y)=1\}.
\]
This is a finite set, and does not contain zero.
Call $H$ a vertex if the linear functionals
$K\mapsto q_K(y)$, for $y\in M(H)$, span $\mathcal H^*$.
Equivalently,
\[
 q_K(y)=0\ (y\in M(H))\quad\Longrightarrow\quad K=0.
 \tag{AC4}
\]
This is the actual vertex definition in the source.  It agrees
with the usual definition as an extreme point of $\mathcal R$.
To see this, if (AC4) holds and $H$ is the midpoint of
$H_1,H_2\in\mathcal R$, all its active inequalities are active
at both endpoints, so $H_1-H_2=0$.
Conversely, if a nonzero $K$ annihilates all active inequalities,
(AC3) and the finitely many remaining strict inequalities show
that $H\pm\varepsilon K\in\mathcal R$ for sufficiently small
$\varepsilon>0$.  Thus $H$ is not extreme.
We will repeatedly use this last perturbation argument.

\subsection{Normal cones and an exact ordered complex}

Let $V$ be the set of vertices of $\mathcal R$.  Define
\[
 \mathcal U=\{\ell\in\mathcal H^*:
            \ell(K)>0\text{ for all }0\ne K\in\mathcal P\},
 \qquad
 C_H=\{\ell:\ell(H)\le\ell(K)\text{ for all }K\in\mathcal R\}.
 \tag{AC5}
\]
The cone $\mathcal U$ is convex and nonempty: the sum of the
matrix traces belongs to it.  Each $C_H$ is closed and convex.

\begin{lemma}\label{ac-normal-cover}
The sets $C_H$, $H\in V$, cover $\mathcal U$.
For every finite nonempty set $E\subset\mathcal U$, only finitely
many of these cones meet $\operatorname{conv}(E)$.
\end{lemma}
\begin{proof}
Fix $\ell\in\mathcal U$.  On the intersection of $\mathcal P$
with the unit sphere, $\ell$ has a strictly positive minimum.
It follows that $\ell(K)\ge c\|K\|$ on $\mathcal P$ for some
$c>0$.  Therefore every sublevel
\[
          \{K\in\mathcal R:\ell(K)\le a\}
\]
is compact.  As $\mathcal R$ is nonempty, $\ell$ attains its
minimum there, and its minimizing face is nonempty and compact.
A nonempty compact convex set in a finite-dimensional real
space has an extreme point: one can successively maximize the
coordinates of a fixed basis, restricting to the maximizing
face after each step, until a singleton remains.
An extreme point of this minimizing face is extreme in
$\mathcal R$, because the value of $\ell$ at each endpoint
of any segment with midpoint in the face must also be minimal.
Hence it belongs to $V$, proving the covering assertion.

We need one consequence of (AC3).  A compact subset
$K\subset\mathcal R$ contains only finitely many vertices.
Cover $K$ by finitely many neighborhoods as in (AC3), and let
$T$ be the union of their finite sets of inequalities.
Every active set $M(H)$ for $H\in K$ is contained in $T$.
For a vertex, this set determines $H$ uniquely by the
equations $q_H(y)=1$.  There are only finitely many subsets
of $T$.

Now fix $H_0\in\mathcal R$, and write
$E=\{\ell_1,\ldots,\ell_t\}$.  If $C_H$ meets
$\operatorname{conv}(E)$, some convex combination
$\ell=\sum_i a_i\ell_i$ is minimized at $H$.  Thus
\[
                 \sum_i a_i\{\ell_i(H)-\ell_i(H_0)\}\le0.
\]
For some $i$ with $a_i>0$, $\ell_i(H)\le\ell_i(H_0)$.
The vertex $H$ therefore lies in the finite union of compact
sublevels defined by these $t$ inequalities.  The preceding
paragraph proves the desired finiteness.
\end{proof}

For $p\ge0$, define $X_p$ to be the set of ordered tuples
$(H_0,\ldots,H_p)$ of vertices such that
\[
                   \mathcal U\cap C_{H_0}\cap\cdots
                            \cap C_{H_p}\ne\varnothing.
                   \tag{AC6}
\]
Repetitions are allowed.  Deleting an entry defines a map
$X_p\to X_{p-1}$.  Put
\[
 C_p=\mathbb Z[X_p],\quad
 \partial(H_0,\ldots,H_p)
      =\sum_{i=0}^p(-1)^i(H_0,\ldots,\widehat H_i,\ldots,H_p),
 \quad
 \epsilon(H)=1.
 \tag{AC7}
\]
Here $\mathbb Z[X]$ denotes the free abelian group on $X$,
with finite support.  The familiar cancellation of pairs of
deletions gives $\partial^2=0$.

We prove the needed exactness directly, including the issue
that the cover consists of closed sets.

\begin{lemma}\label{ac-convex-nerve}
Let $K$ be a nonempty compact convex subset of a
finite-dimensional real vector space, covered by finitely many
closed convex sets $A_i$.  Form the ordered tuple complex of
the nerve: a tuple is allowed precisely when
$K\cap\bigcap_j A_{i_j}\ne\varnothing$.  Its augmented integral
chain complex, with repetitions and the boundary (AC7), is exact.
\end{lemma}
\begin{proof}
Replace $A_i$ by $K\cap A_i$ and discard empty sets.
We induct on the number of subsets of the finite index set
whose corresponding intersection is nonempty, counting the
empty subset as well.

If all the $A_i$ have a common point, every tuple of indices
is allowed.  Fix an index $a$.  Prepending $a$ to a tuple
defines a contraction: for $p\ge0$,
\[
       \partial(a,i_0,\ldots,i_p)
          =(i_0,\ldots,i_p)-(a,\partial(i_0,\ldots,i_p)),
\]
and in degree $-1$ send $1$ to $(a)$.  This proves exactness.

Otherwise choose a minimal collection of the $A_i$ with
empty intersection.  Write it as $\{A_i\}\cup\{A_j:j\in J\}$.
Both $A_i$ and $B=\bigcap_{j\in J}A_j$ are nonempty compact
convex sets, and they are disjoint.
They admit a strict separating linear functional $f$:
\[
                         f(A_i)<u<f(B)
                         \tag{AC8}
\]
for some real number $u$.
For completeness, take a pair of nearest points $a\in A_i$,
$b\in B$.  Their distance is positive by compactness.
With $f(x)=\langle b-a,x\rangle$, differentiating the squared
distance along segments in the two convex sets gives
$f(x)\le f(a)<f(b)\le f(y)$ for $x\in A_i$, $y\in B$.
Any number between $f(a)$ and $f(b)$ gives (AC8).

Set $K_-=K\cap\{f\le u\}$, $K_+=K\cap\{f\ge u\}$,
and $K_0=K\cap\{f=u\}$.  They are nonempty compact convex
sets, covered by the corresponding intersections with $A_i$.
All three nerves have strictly fewer faces: $K_-$ and $K_0$
lose the face $J$, and $K_+$ loses the vertex $i$.
Their augmented complexes are therefore exact by induction.

The nerve on $K$ is the union of the nerves on $K_-$ and
$K_+$.  Their intersection is the nerve on $K_0$.
Indeed, witnesses on both sides for the same tuple can be
joined by a segment lying in every set of that tuple, and
this segment meets the separating hyperplane.
These statements remain true for ordered tuples with repetitions.

Here is the chain argument for this union.  Decompose a cycle
$z$ as $z_-+z_+$, using chains in the two subcomplexes.
Then $\partial z_-=-\partial z_+$ is supported in their
intersection.  Exactness there supplies a chain $w$ with
$\partial w=\partial z_-$.  The chains $z_--w$ and $z_++w$
are cycles in their respective subcomplexes, hence boundaries.
Their sum is $z$.  In degree zero the same argument uses
surjectivity of the augmentation of the intersection; in
degree $-1$ surjectivity follows from nonemptiness.
This completes the induction.
\end{proof}

\begin{theorem}\label{ac-exact}
The augmented sequence
\[
       \cdots\longrightarrow\mathbb Z[X_2]
       \longrightarrow\mathbb Z[X_1]
       \longrightarrow\mathbb Z[X_0]
       \overset{\epsilon}{\longrightarrow}\mathbb Z
       \longrightarrow0
\]
is an exact sequence of permutation $\mathbb Z\Gamma$-modules.
Every column of the boundary
$\mathbb Z[X_{p+1}]\to\mathbb Z[X_p]$ has integral
$\ell^1$-norm at most $p+2$.
\end{theorem}
\begin{proof}
The action (AC2), and its contragredient action on linear
functionals, preserve (AC6).  Deletions and augmentation are
equivariant, so the modules are permutation modules.
The boundary estimate counts its $p+2$ signed terms;
coinciding faces can only decrease the $\ell^1$-norm.

A cycle has finite support.  Choose a witness in $\mathcal U$
for each tuple in its support, and choose an additional point
of $\mathcal U$ if necessary.  Their convex hull $K$ is
nonempty, compact, and contained in $\mathcal U$.
By Lemma~\ref{ac-normal-cover}, only finitely many $C_H$
meet $K$, and they cover $K$.  Apply
Lemma~\ref{ac-convex-nerve} to their intersections with $K$.
It supplies a filling using only tuples already allowed
by (AC6).  The same covering lemma ensures that $X_0$
is nonempty, so augmentation is surjective.
\end{proof}

\subsection{Common minimal vectors and determinant bounds}

\begin{lemma}\label{ac-common-minima}
For every tuple $(H_0,\ldots,H_p)\in X_p$, the common minimal
vectors
\[
                 \bigcap_{j=0}^p M(H_j)
\]
span $F^N$ over $F$.
\end{lemma}
\begin{proof}
Let $\ell\in\mathcal U$ witness (AC6), and put
$\overline H=(p+1)^{-1}\sum_jH_j$.
Each $H_j$ minimizes $\ell$ on $\mathcal R$, so
$\overline H$ does too.  Since every $q_{H_j}(y)\ge1$,
\[
                  M(\overline H)=\bigcap_jM(H_j).
\]
Suppose this set does not span $F^N$.  There exists a nonzero
$F$-linear functional $a:F^N\to F$ vanishing on it.
At each place let $K_v$ be the positive semidefinite rank-one
form whose value at $z$ is $|\sigma_v(a)z|^2$.
This defines a nonzero $K\in\mathcal P$ with
$q_K(y)=0$ for every $y\in M(\overline H)$.
The perturbation argument following (AC4) gives
$\overline H-\varepsilon K\in\mathcal R$ for some
$\varepsilon>0$.  But
$\ell(\overline H-\varepsilon K)<\ell(\overline H)$,
a contradiction.
\end{proof}

Consequently we may choose $N$ common minimal vectors as the
columns of an integral matrix $A\in M_N(\mathcal O_F)$ with
$\det A\ne0$.  Such columns are an $F$-basis, and are not
claimed to be an $\mathcal O_F$-basis of $\Lambda$.
Define the integers
\[
       I_0=D^N d^d,\qquad R_0=N^N I_0^2,\qquad
       I(A)=|\operatorname{Norm}_{F/\mathbb Q}(\det A)|.
       \tag{AC9}
\]

\begin{lemma}\label{ac-det}
For every $H\in\mathcal R$,
\[
                  I_0\prod_v\det H_v\ge1.
                  \tag{AC10}
\]
If the columns of $A$ are minimal for $H$, put
$H'_v=\sigma_v(A)^*H_v\sigma_v(A)$.  Then
\[
 \begin{split}
 &1\le I(A)\le I_0,\\
 &H'_v>0,\qquad (H'_v)_{aa}\le1,\qquad
                I_0^{-1}\le\det H'_v\le1
                   \quad\text{for every }v,a.
 \end{split}
 \tag{AC11}
\]
\end{lemma}
\begin{proof}
Let $G_H$ be the real Gram matrix in the integral coordinates
used above.  Its determinant is
\[
       \det G_H=2^{-d}D^N\left(\prod_v\det H_v\right)^2.
       \tag{AC12}
\]
Indeed the squared determinant of the lattice embedding is
$2^{-d}D^N$.  For a positive Hermitian matrix, its quadratic
form on the underlying real space has determinant equal to
the square of its complex determinant: unitary diagonalization
reduces this to the fact that each eigenvalue occurs twice.
This proves (AC12).

Write $P=\prod_v\det H_v>0$.
Lemma~\ref{ac-positive-gram} gives
$P^2\ge 2^dD^{-N}d^{-d}$, and hence
\[
                       (I_0P)^2\ge D^N(2d)^d\ge1.
\]
This proves (AC10); it also shows that our intermediate
determinant estimate is stronger than the one required.

For each column $a$ of $A$, minimality says
$\sum_v(H'_v)_{aa}=1$.  Every term is positive, so
$(H'_v)_{aa}\le1$.  Hadamard's inequality gives
$0<\det H'_v\le1$.  This inequality follows, for example,
by writing $H'_v$ as a Gram matrix: successive orthogonal
projections do not increase the lengths of its Gram vectors.
The determinant transformation formula gives
\[
                  \prod_v\det H'_v=I(A)P.
                  \tag{AC13}
\]
Here each complex embedding is represented once, so
$\prod_v|\sigma_v(\det A)|^2=I(A)$.
The left side of (AC13) is at most one, whence
$I(A)\le P^{-1}\le I_0$.
Each factor on that left side is at least their product,
because all the other factors are at most one.  Since
$I(A)\ge1$, this proves
$\det H'_v\ge I(A)P\ge I_0^{-1}$.
\end{proof}

\subsection{Integral matrix reduction without a principal ideal hypothesis}

\begin{lemma}\label{ac-scaled-inverse}
For $A\in M_N(\mathcal O_F)$ with $\det A\ne0$, the matrix
\[
                           B_A=I(A)A^{-1}
                           \tag{AC14}
\]
belongs to $M_N(\mathcal O_F)$.
For every integer $J\ge1$, there is a set $\mathscr A_J$
of at most $J^{1+dN}$ integral matrices such that every
$A$ with $0<I(A)\le J$ satisfies
$gA\in\mathscr A_J$ for some $g\in\Gamma$.
\end{lemma}
\begin{proof}
If $0\ne\alpha\in\mathcal O_F$, the group
$\mathcal O_F/\alpha\mathcal O_F$ has order
$|\operatorname{Norm}_{F/\mathbb Q}(\alpha)|$.
Indeed multiplication by $\alpha$ on an integral basis has
that determinant, and the index of an integral full-rank
matrix is the absolute value of its determinant, as follows
from integer elementary row and column reduction.
Lagrange's theorem in this finite additive group therefore
gives
\[
          |\operatorname{Norm}_{F/\mathbb Q}(\alpha)|
                      \in\alpha\mathcal O_F.
\]
For $\alpha=\det A$, the adjugate formula now proves (AC14).
In particular $AB_A=B_AA=I(A)I_N$.

Fix a positive integer $I$, and consider only matrices with
$I(A)=I$.  Assign to $A$ the residue of $B_A$ in
$M_N(\mathcal O_F/I\mathcal O_F)$.
If $A,A'$ have the same residue, write
$B_A-B_{A'}=IX$ with $X\in M_N(\mathcal O_F)$.
Then
\[
       A'A^{-1}=I_N+A'X,\qquad
       AA'^{-1}=I_N-AX.
       \tag{AC15}
\]
Both matrices on the right are integral and are inverse
to each other.  Thus $g=A'A^{-1}$ belongs to $\Gamma$
and satisfies $gA=A'$.

There are exactly $I^{nN^2}=I^{dN}$ possible residues,
because the additive rank of $\mathcal O_F$ is $n$.
Choose one representative matrix for each residue that
actually occurs.  Equal residues imply equal left
$\Gamma$-orbits by (AC15), so these representatives meet
every such orbit.  They need not be distinct orbit
representatives, which only improves the upper bound.
Summing over $1\le I\le J$ gives at most
\[
                         \sum_{I=1}^J I^{dN}
                             \le J^{1+dN}
\]
matrices.  No ideal in this argument is required to be
principal except the explicitly displayed ideal
$\alpha\mathcal O_F$.
\end{proof}

\subsection{Counting vertices with fixed minimal columns}

For an integral matrix $A$ with nonzero determinant, define
\[
 \begin{split}
 T(A)=\{y\in\Lambda:\;&q_H(y)=1
       \text{ for some }H\in\mathcal R\\
       &\text{for which every column of }A\text{ is minimal}\},
 \end{split}
\]
and let $V(A)$ be the set of vertices for which every column
of $A$ is minimal.  These sets are allowed to be empty.

\begin{lemma}\label{ac-cramer}
Let $G$ be a positive definite Hermitian $N$ by $N$ matrix
such that $G_{aa}\le1$ for every $a$.  If $c^*Gc\le1$, then
\[
                          (\det G)|c_a|\le N^N
                          \quad(1\le a\le N).
                          \tag{AC16}
\]
\end{lemma}
\begin{proof}
Cauchy--Schwarz for the inner product defined by $G$ gives
$|G_{ij}|\le1$ and $|(Gc)_i|\le1$.
Cramer's rule says that $(\det G)c_a$ is the determinant
of the matrix obtained by replacing column $a$ of $G$
by $Gc$.  Every entry of this matrix has absolute value
at most one.  Its determinant is a sum of $N!$ signed
products, so its absolute value is at most
$N!\le N^N$.
\end{proof}

\begin{lemma}\label{ac-short-count}
One has
\[
             |T(A)|\le(2R_0+1)^d,\qquad
             |V(A)|\le\bigl((2R_0+1)^d+1\bigr)^{dN}.
             \tag{AC17}
\]
In particular both sets are finite.
\end{lemma}
\begin{proof}
Suppose $y\in T(A)$, and choose the form $H$ in its definition.
Write $I=I(A)$, $B=B_A$, and $z=By$.
Since $Az=Iy$, and $H'_v=\sigma_v(A)^*H_v\sigma_v(A)$,
\[
              \sum_v \sigma_v(z)^*H'_v\sigma_v(z)=I^2.
\]
Every term is nonnegative.  For each place apply
Lemma~\ref{ac-cramer} to $G=H'_v$ and
$c=I^{-1}\sigma_v(z)$.  Using (AC11) gives
\[
       |\sigma_v(z_a)|
          \le I N^N I_0
          \le N^N I_0^2=R_0
                  \qquad(v,a).
                  \tag{AC18}
\]

For any integer $R\ge0$, there are at most $(2R+1)^n$
algebraic integers $a\in\mathcal O_F$ satisfying
$|\sigma_v(a)|\le R$ for all $v$.
To prove this, reduce modulo $(2R+1)\mathcal O_F$.
If two such integers $a,b$ have the same residue, write
$a-b=(2R+1)c$ with $c\in\mathcal O_F$.  Then
\[
                    |\sigma_v(c)|\le\frac{2R}{2R+1}<1
\]
at every place.  The norm of $c$ is an integer of absolute
value less than one, so $c=0$.  Reduction is injective on
the set being counted, and the quotient has
$(2R+1)^n$ elements.  Applying this with $R=R_0$ in each
of the $N$ coordinates bounds the possible $z$ by
$(2R_0+1)^d$.  The map $y\mapsto B_Ay$ is injective
because $B_A$ is invertible over $F$.  This proves the
first assertion in (AC17).

If $H\in V(A)$, its active functionals span
$\mathcal H^*$ by definition.  Choose at most
$\dim_{\mathbb R}\mathcal H\le dN$ of its minimal vectors
whose functionals still span this dual space.
All of them belong to $T(A)$, and the equations
$q_K(y)=1$ for this chosen set uniquely determine $K=H$.
Choose one such determining subset for each vertex.
Two vertices assigned the same subset must be equal.
A set of $t$ elements has at most $(t+1)^r$ subsets of
size at most $r$: order the elements, list each subset
in order, and pad its list to length $r$ with an extra
symbol.  Taking $t=|T(A)|$ and $r=dN$ proves the
second assertion.
\end{proof}

\begin{theorem}\label{ac-orbit}
For every $p\ge0$, the ordered tuple set $X_p$ has finitely
many $\Gamma$-orbits, and
\[
 \begin{split}
 \#(\Gamma\backslash X_p)
   &\le I_0^{\,1+dN}
       \bigl((2R_0+1)^d+1\bigr)^{dN(p+1)}\\
   &\le (2Dd)^{12(p+1)d^3N}.
 \end{split}
 \tag{AC19}
\]
\end{theorem}
\begin{proof}
For a tuple in $X_p$, choose the matrix $A$ of common
minimal columns furnished by Lemma~\ref{ac-common-minima}.
By Lemma~\ref{ac-det}, $I(A)\le I_0$.
Choose $g\in\Gamma$ such that $gA$ belongs to the set
$\mathscr A_{I_0}$ in Lemma~\ref{ac-scaled-inverse}.
Because $q_{gH_j}(gy)=q_{H_j}(y)$, every vertex of the
translated tuple belongs to $V(gA)$.
Every orbit thus meets
\[
              \bigcup_{A'\in\mathscr A_{I_0}}V(A')^{p+1}.
\]
Only the tuples satisfying (AC6) need be included, so
counting all displayed tuples is an upper bound.
Lemma~\ref{ac-short-count} proves the first inequality.

Here are the numerical estimates for the second inequality.
Put $b=2Dd\ge2$.  Since $1\le N\le d$,
\[
 \begin{aligned}
 I_0=D^N d^d&\le b^{2d},&
 N^N&\le b^d,&
 R_0&\le b^{5d},\\
 2R_0+1&\le b^{5d+2},&
 (2R_0+1)^d+1&\le b^{(5d+2)d+1}.
 \end{aligned}
\]
For the fourth inequality use $R_0\ge1$ and
$2R_0+1\le3R_0\le b^2R_0$; for the fifth use $2\le b$.
The exponent resulting from the first line of (AC19)
is therefore at most
\[
       2d+2d^2N+(5d^3N+2d^2N+dN)(p+1).
\]
Each of its five terms is bounded respectively by
$2$, $2$, $5$, $2$, and $1$ times $(p+1)d^3N$.
Their sum is at most $12(p+1)d^3N$.
\end{proof}

\subsection{Stabilizers, hypotheses, and the resulting homology estimate}

\begin{proposition}\label{ac-stabilizer}
Every stabilizer of a tuple in $X_p$ is finite and has
order at most $3^{nN^2}=3^{dN}$.
\end{proposition}
\begin{proof}
All vertices are positive definite by
Proposition~\ref{ac-polyhedron}.  At the complex places,
the stabilizer of one such tuple is contained in the
product of the unitary stabilizers of its first entry.
This is a compact subset of the product of complex
matrix spaces.  Matrices over $\mathcal O_F$ form a
closed discrete subset of that real vector space,
so their intersection with this compact set is finite.

For the quantitative assertion apply
Lemma~\ref{ht-congruence-three} to
$M_N(\mathcal O_F)$, whose additive rank is $nN^2$.
That lemma proves that reduction modulo $3$ is injective
on every finite subgroup of its unit group.  The residue
ring has $3^{nN^2}$ elements.
\end{proof}

Theorem~\ref{ac-exact}, Theorem~\ref{ac-orbit}, and
Proposition~\ref{ac-stabilizer} supply all the hypotheses
of Proposition~\ref{ht-permutation} for this arithmetic
group.  In degree $m\ge1$ one may take
\[
       A=(2Dd)^{12(m+1)d^3N},\qquad
       s=3^{dN},\qquad B=m+2.
\]
In particular the geometric input of
Proposition~\ref{ht-arithmetic-numerical} is now proved.
For $E\ge1$, $1\le N\le4E+6$, and $0\le m\le2E+1$,
its established numerical argument gives the unconditional
inequality
\[
 \log\#H_m(\operatorname{GL}_N(\mathcal O_F),\mathbb Z)_{\rm tors}
       \le\exp\bigl((2En)^{13}\log(2D)\bigr).
 \tag{AC20}
\]
No class-number estimate is imported in the proof of
(AC20).  The remaining passage from this homology estimate
to the strict Selmer group in the submitted proof is
a different assertion.

\paragraph{Boundary cases and normalization checks.}
The source also allows $N=0$ in some intermediate statements.
In that case $\Lambda=0$, $\mathcal H=0$, $\mathcal R$
has one point, and $\Gamma$ is trivial.  The ordered
complex has a single generator in every degree and
is contracted by prepending its unique vertex.
Every orbit and stabilizer has cardinality one.
With $0^0=1$, the unsimplified first bound in (AC19)
also equals one.  The simplified bound involving $d$
is only asserted above for $N\ge1$.

The normalization in (AC12) can be checked already at
$F=\mathbb Q(i)$, $N=1$.  For $H=(t)$ and the integral
basis $1,i$, one has $G_H=tI_2$, $D=4$, and
$2^2\det G_H=4t^2=D(\det H)^2$.
Removing the factor $2^{-d}$ would give a false identity.
This check concerns the convention of using one embedding
per conjugate pair; the source includes the correct factor.

Positive definiteness is also essential to the stabilizer
argument.  On $\mathcal O_F^2$, the semidefinite form
$|x_1|^2$ at every complex place is preserved by all
matrices
\[
                      \begin{pmatrix}1&0\\a&1\end{pmatrix},
                         \qquad a\in\mathcal O_F.
\]
Its arithmetic stabilizer is infinite.  It is not a
counterexample to the proposition: the integral vector
$(0,1)$ has value zero, so this form is excluded from
$\mathcal R$ by (AC1).  Proposition~\ref{ac-polyhedron}
proves exactly the positivity needed to exclude this
possibility.

Finally, the construction uses the ordered tuples (AC6),
with repetitions and the unsigned action on their
generators.  Replacing them by oriented simplices without
accounting for orientation characters would change the
permutation-module argument.  The proof above makes no
such replacement.  No counterexample satisfying the actual
hypotheses was found.




